\section{International Legal Framework}

\lead{The governance of outer space rests upon a series of multilateral treaties negotiated principally under the auspices of the United Nations Committee on the Peaceful Uses of Outer Space (COPUOS).}

These instruments form the backbone of international space law. The following summarises the key instruments and their relevance to the topic of privatisation and commercial colonisation.

\subsection{The Outer Space Treaty (1967)}

\emph{Full title:} Treaty on Principles Governing the Activities of States in the Exploration and Use of Outer Space, including the Moon and Other Celestial Bodies.\footnote{United Nations Office for Outer Space Affairs. (n.d.). \emph{Outer Space Treaty}. \url{https://www.unoosa.org/oosa/en/ourwork/spacelaw/treaties/outerspacetreaty.html}}

This is the foundational instrument of international space law. Its key principles include:

\begin{itemize}
  \item \textbf{Non-appropriation:} Outer space, including the Moon and other celestial bodies, is not subject to national appropriation by claim of sovereignty, by means of use or occupation, or by any other means (Article II).
  \item \textbf{Freedom of exploration:} Outer space shall be free for exploration and use by all States (Article I).
  \item \textbf{Peaceful use:} Celestial bodies shall be used exclusively for peaceful purposes (Article IV).
  \item \textbf{State responsibility (Article VI):} States bear international responsibility for national activities in outer space, including those of non-governmental (private) entities. Private activities require authorisation and continuing supervision by the appropriate State.
  \item \textbf{Liability (Article VII):} A State that launches or procures the launch of an object into outer space is internationally liable for damage caused by that object.
\end{itemize}

\emph{\textbf{Article II, Outer Space Treaty}}\emph{: "Outer space, including the Moon and other celestial bodies, is not subject to national appropriation by claim of sovereignty, by means of use or occupation, or by any other means."}

\subsection{The Rescue Agreement (1968)}

The Agreement on the Rescue of Astronauts, the Return of Astronauts and the Return of Objects Launched into Outer Space requires States to render all possible assistance to astronauts in distress and to return space objects found outside the launching State's territory.

\subsection{The Liability Convention (1972)}

The Convention on International Liability for Damage Caused by Space Objects elaborates on Article VII of the Outer Space Treaty. It establishes an absolute liability regime for damage caused on the surface of the Earth and a fault-based regime for damage caused in outer space.\footnote{European Space Law. (n.d.). \emph{International}. \url{https://www.european-space-law.com/International.html}}

\subsection{The Registration Convention (1976)}

The Convention on Registration of Objects Launched into Outer Space requires launching States to register space objects with the United Nations, facilitating the identification and tracking of objects in orbit.

\subsection{The Moon Agreement (1979)}

The Agreement Governing the Activities of States on the Moon and Other Celestial Bodies declares that the Moon and its natural resources are the “common heritage of mankind” and calls for the establishment of an international regime to govern resource exploitation.\footnote{Nuclear Threat Initiative. (n.d.). \emph{Moon Agreement}. \url{https://www.nti.org/education-center/treaties-and-regimes/agreement-governing-activities-states-moon-and-other-celestial-bodies-moon-agreement/}} Crucially, few major spacefaring States (including the United States) have ratified this agreement, limiting its practical effect.\footnote{Space Legal Issues. (n.d.). \emph{The 1979 Moon Agreement}. \url{https://www.spacelegalissues.com/the-1979-moon-agreement/}}

\subsection{The Artemis Accords (2020)}

The Artemis Accords are a U.S.-led set of bilateral principles for cooperative civil space exploration. They are not a treaty but rather a political framework that complements the Outer Space Treaty.\footnote{U.S. Department of State. (n.d.). \emph{Artemis Accords}. \url{https://www.state.gov/bureau-of-oceans-and-international-environmental-and-scientific-affairs/artemis-accords}} Key features include:

\begin{itemize}
  \item Affirmation that space resource extraction does not inherently constitute national appropriation.
  \item Establishment of “safety zones” around operations to avoid harmful interference.
  \item Commitments to transparency, interoperability, and registration of space objects.
\end{itemize}

\subsection{Friction Between U.S. Domestic Law and International Principles}

The U.S. domestic statutes recognising private property rights in extracted space resources (most notably the SPACE Act of 2015) represent a deliberate interpretation that the non-appropriation principle applies to sovereignty claims over celestial bodies themselves, but does not preclude ownership of resources once extracted, analogous to fishing rights in the high seas.\footnote{Reddie \& Grose. (2024, October 9). \emph{Outer space resource acquisition – A legal perspective}. \url{https://www.reddie.co.uk/2024/10/09/outer-space-resource-acquisition-a-legal-perspective/}} Critics argue that this interpretation stretches the Outer Space Treaty beyond its intended scope and risks undermining the multilateral framework.\footnote{Emory International Law Review. (n.d.). \emph{Commercial Space Launch Competitiveness Act of 2015 and the exploitation of resources in outer space}. \url{https://scholarlycommons.law.emory.edu/cgi/viewcontent.cgi?article=1218&context=eilr}} The Artemis Accords seek to build international consensus around this interpretation, but they remain controversial among some States.
